% ----------------------
\subsection{Section 39 (5 January 1831)}
% ----------------------
	\label{DC39}
	
	% -----
	\subsubsection{Historical Background}
	% -----
		\label{DC39-HB}
		
		% --
		\paragraph{JSP}
		% --
		
			``JS dictated this revelation in Fayette, New York, for James Covel, a Protestant minister, three days after the Church of Christ's third conference. When John Whitmer recorded this text in Revelation Book 1 a few months later, he wrote that Covel `covenanted with the Lord that he would obey any commandment that the Lord would give through his servent Joseph.'
			
			``The identity of the revelation's recipient is not known with certainty. Two individuals living in New York at the time fit the general description, and no source definitively identifies either man as the recipient. The earliest extant manuscript copy of the revelation, featured below, provides only a given name. The first printed version in 1833 expanded `James' to `James (C.,)' with no additional information, and in 1835 the name was given as `James Covill' in the Doctrine and Covenants. JS's history also uses this spelling because its editors relied on the 1835 edition of the Doctrine and Covenants for the revelation text. The history adds that Covill `had been a baptist minister for about forty years.' James Covill, a Baptist minister from Ellery, New York, who in 181 was over seventy years old, fits this description, but he lived on the far western edge of the state, more than one hundred fifty miles away. JS and Sidney Rigdon could have met Covill on their way to Ohio at the end of January, but according to this earliest copy of the revelation, it was `given at Fayette' on 5 January 1831.
			
			The recipient of the revelation was much more likely James Covel, a Methodist elder from Canadice, New York. The index to Revelation Book 1 describes the recipient as `a Methodist Priest,' not a Baptist. James Covel lived about twenty miles southwest of Canandaigua, New York, and had been associated with the Methodist church for nearly forty years. He may have heard JS or Sidney Rigdon preaching in the Canandaigua area. After JS and several others preached `with great power' in Ezra Thayer's barn near Canandaigua in October 1830, they were invited to preach in Canandaigua. `They had promised that we should meet in the Methodist Meeting house,' Thayer later wrote, `but the Trustees could not agree.' As president of the regional Methodist conference, Covel was likely aware of the request. In December a Mormon preacher, probably JS or Rigdon, `delivered a discourse in the Town House [in Canandaigua] to an assembly of two or three hundred people.' Covel may have attended the December meeting and then traveled to Fayette, where the revelation was dictated.
			
			Within a day after JS dictated this revelation, Covel departed from Fayette without explanation, leaving JS and Rigdon to wonder why he did not follow the commandment. A revelation on 6 January explained `why he obeyed not the word.'''%
				\footnote{%
					% \label{}%
					See \hyperref[DC40]{DC 40}.
					}
			
	% -----
	\subsubsection{Versions}
	% -----
		\label{DC39-V}
		
		See the \href{https://drive.google.com/file/d/1goAvNz8jS4R8slYfMG_db55Ty2z_vmgK/view?usp=sharing}{sections comparison} document.
	
		\begin{compactdesc}
			\item[JSP] \hyperref[EMS-1831-JS-5-JAN-1831]{5 January 1831}
			\item[BC] Chapter 41, 85--87
			\item[DC 1835] Section 59, 187--188
			\item[TS] 4:23 (15 October 1843), 353--354
			\item[MS] 5:3 (August 1844), 35
			\item[DC 1844] Section 59, 282--285
		\end{compactdesc}

	% -----
	\subsubsection{Section 39:1} % *DC39-1 
	% -----
		\label{DC39-1}

		\hyperref[HEARKEN]{HEARKEN} and listen to the voice of him who is from all eternity to all eternity,%
			\footnote{%
				% \label{}%
				For the phrase ``eternity to eternity,'' see \hyperref[ETERNITY-PHRASES]{here}.
				}
		the Great I AM,%
			\footnote{%
				% \label{}%
				See \hyperref[DC29-1]{DC 29:1} and notes there.
				}
		even Jesus Christ---

		% \begin{framed}
		% \scriptsize
			% \begin{compactdesc}
				% \item[JSP] 
				% % \item[BC] 
				% % \item[]
			% \end{compactdesc}
		% \end{framed}

	% -----
	\subsubsection{Section 39:2} % *DC39-2 
	% -----
		\label{DC39-2}

		The light and the life of the world;%
			\footnote{%
				% \label{}%
				See \hyperref[JOHN-1-4]{John 1:4} and the phrase ``\hyperref[LIGHT-PHRASES]{light and life}.'' Compare also the beginning of \hyperref[DC34]{DC 34} and \hyperref[DC45]{DC 45}.
				}
		a light which shineth in darkness and the darkness comprehendeth it not;%
			\footnote{%
				% \label{}%
				See \hyperref[JOHN-1-5]{John 1:5}.
				}

		% \begin{framed}
		% \scriptsize
			% \begin{compactdesc}
				% \item[JSP] 
				% % \item[BC] 
				% % \item[]
			% \end{compactdesc}
		% \end{framed}

	% -----
	\subsubsection{Section 39:3} % *DC39-3 
	% -----
		\label{DC39-3}

		The same which came in the meridian%
			\footnote{%
				% \label{}%
				This word appears six times in the scriptures, none in the Bible---here, \hyperref[DC20-26]{DC 20:26}, \hyperref[MOSES-5-57]{Moses 5:57}, \hyperref[MOSES-6-57]{6:57}, \hyperref[MOSES-6-62]{6:62}, and \hyperref[MOSES-7-46]{7:46}. From the 1828 dictionary:
					\begin{compactitem}
						\item[1]  In astronomy and geography, a great circle supposed to be drawn or to pass through the poles of the earth, and the zenith and nadir of any given place, intersecting the equator at right angles, and dividing the hemisphere into eastern and western. Every place on the globe has its \textit{meridian} and when the sun arrives at this circle, it is mid-day or noon. whence the name. This circle may be considered to be drawn on the surface of the earth, or it may be considered as a circle in the heavens coinciding with that on the earth.
						\item[2] Mid-day, noon.
						\item[3] The highest point; as the \textit{meridian} of life; the \textit{meridian} of power or of glory.
						\item[4] The particular place or state, with regard to local circumstances or things that distinguish it from others. We say, a book is adapted to the \textit{meridian} of France or Italy; a measure is adapted to the \textit{meridian} of London or Washington.
					\end{compactitem}
				}
		of time unto mine own, and mine own received me not;

		% \begin{framed}
		% \scriptsize
			% \begin{compactdesc}
				% \item[JSP] 
				% % \item[BC] 
				% % \item[]
			% \end{compactdesc}
		% \end{framed}

	% -----
	\subsubsection{Section 39:4} % *DC39-4 
	% -----
		\label{DC39-4}

		But to as many as received me, gave I power to become my sons; and even so will I give unto as many as will receive me, power to become my sons.

		% \begin{framed}
		% \scriptsize
			% \begin{compactdesc}
				% \item[JSP] 
				% % \item[BC] 
				% % \item[]
			% \end{compactdesc}
		% \end{framed}

	% -----
	\subsubsection{Section 39:5} % *DC39-5 
	% -----
		\label{DC39-5}

		And verily, verily, I say unto you, he that receiveth my gospel receiveth me; and he that receiveth not my gospel receiveth not me.

		% \begin{framed}
		% \scriptsize
			% \begin{compactdesc}
				% \item[JSP] 
				% % \item[BC] 
				% % \item[]
			% \end{compactdesc}
		% \end{framed}

	% -----
	\subsubsection{Section 39:6} % *DC39-6 
	% -----
		\label{DC39-6}

		And this is my gospel---repentance and baptism by water, and then cometh the baptism of fire and the Holy Ghost, even the \hyperref[COMFORTER]{Comforter}, which showeth all things, and teacheth the \hyperref[PEACEABLE]{peaceable} things of the kingdom.

		% \begin{framed}
		% \scriptsize
			% \begin{compactdesc}
				% \item[JSP] 
				% % \item[BC] 
				% % \item[]
			% \end{compactdesc}
		% \end{framed}

	% -----
	\subsubsection{Section 39:7} % *DC39-7 
	% -----
		\label{DC39-7}

		And now, behold, I say unto you, my servant James, I have looked upon thy works and I know thee.

		% \begin{framed}
		% \scriptsize
			% \begin{compactdesc}
				% \item[JSP] 
				% % \item[BC] 
				% % \item[]
			% \end{compactdesc}
		% \end{framed}

	% -----
	\subsubsection{Section 39:8} % *DC39-8 
	% -----
		\label{DC39-8}

		And verily I say unto thee, thine heart is now right before me at this time; and, behold, I have bestowed great blessings upon thy head;

		% \begin{framed}
		% \scriptsize
			% \begin{compactdesc}
				% \item[JSP] 
				% % \item[BC] 
				% % \item[]
			% \end{compactdesc}
		% \end{framed}

	% -----
	\subsubsection{Section 39:9} % *DC39-9 
	% -----
		\label{DC39-9}

		Nevertheless, thou hast seen great sorrow, for thou hast rejected me many times because of pride and the cares of the world.

		% \begin{framed}
		% \scriptsize
			% \begin{compactdesc}
				% \item[JSP] 
				% % \item[BC] 
				% % \item[]
			% \end{compactdesc}
		% \end{framed}

	% -----
	\subsubsection{Section 39:10} % *DC39-10 
	% -----
		\label{DC39-10}

		But, behold, the days of thy deliverance are come, if thou wilt hearken to my voice, which saith unto thee: Arise and be baptized, and wash away your sins, calling on my name, and you shall receive my Spirit, and a blessing so great as you never have known.

		% \begin{framed}
		% \scriptsize
			% \begin{compactdesc}
				% \item[JSP] 
				% % \item[BC] 
				% % \item[]
			% \end{compactdesc}
		% \end{framed}

	% -----
	\subsubsection{Section 39:11} % *DC39-11 
	% -----
		\label{DC39-11}

		And if thou do this, I have prepared thee for a greater work.  Thou shalt preach the \hyperref[FULNESS]{fulness} of my gospel, which I have sent forth in these last days, the covenant which I have sent forth to recover my people, which are of the house of Israel.

		% \begin{framed}
		% \scriptsize
			% \begin{compactdesc}
				% \item[JSP] 
				% % \item[BC] 
				% % \item[]
			% \end{compactdesc}
		% \end{framed}

	% -----
	\subsubsection{Section 39:12} % *DC39-12 
	% -----
		\label{DC39-12}

		And it shall come to pass that power shall rest upon thee; thou shalt have great faith, and I will be with thee and go before thy face.

		% \begin{framed}
		% \scriptsize
			% \begin{compactdesc}
				% \item[JSP] 
				% % \item[BC] 
				% % \item[]
			% \end{compactdesc}
		% \end{framed}

	% -----
	\subsubsection{Section 39:13} % *DC39-13 
	% -----
		\label{DC39-13}

		Thou art called to labor in my vineyard, and to build up my church, and to bring forth Zion, that it may rejoice upon the hills and flourish.

		% \begin{framed}
		% \scriptsize
			% \begin{compactdesc}
				% \item[JSP] 
				% % \item[BC] 
				% % \item[]
			% \end{compactdesc}
		% \end{framed}

	% -----
	\subsubsection{Section 39:14} % *DC39-14 
	% -----
		\label{DC39-14}

		Behold, verily, verily, I say unto thee, thou art not called to go into the eastern countries, but thou art called to go to the Ohio.

		% \begin{framed}
		% \scriptsize
			% \begin{compactdesc}
				% \item[JSP] 
				% % \item[BC] 
				% % \item[]
			% \end{compactdesc}
		% \end{framed}

	% -----
	\subsubsection{Section 39:15} % *DC39-15 
	% -----
		\label{DC39-15}

		And inasmuch as my people shall assemble themselves at the Ohio, I have kept in store a blessing such as is not known among the children of men, and it shall be poured forth upon their heads.  And from thence men shall go forth into all nations.

		% \begin{framed}
		% \scriptsize
			% \begin{compactdesc}
				% \item[JSP] 
				% % \item[BC] 
				% % \item[]
			% \end{compactdesc}
		% \end{framed}

	% -----
	\subsubsection{Section 39:16} % *DC39-16 
	% -----
		\label{DC39-16}

		Behold, verily, verily, I say unto you, that the people in Ohio call upon me in much faith, thinking I will stay my hand in judgment upon the nations, but I cannot deny my word.

		% \begin{framed}
		% \scriptsize
			% \begin{compactdesc}
				% \item[JSP] 
				% % \item[BC] 
				% % \item[]
			% \end{compactdesc}
		% \end{framed}

	% -----
	\subsubsection{Section 39:17} % *DC39-17 
	% -----
		\label{DC39-17}

		Wherefore lay to with your might and call faithful laborers into my vineyard, that it may be pruned for the last time.

		% \begin{framed}
		% \scriptsize
			% \begin{compactdesc}
				% \item[JSP] 
				% % \item[BC] 
				% % \item[]
			% \end{compactdesc}
		% \end{framed}

	% -----
	\subsubsection{Section 39:18} % *DC39-18 
	% -----
		\label{DC39-18}

		And inasmuch as they do repent and receive the \hyperref[FULNESS]{fulness} of my gospel, and become \hyperref[SANCTIFICATION]{sanctified}, I will stay mine hand in judgment.

		% \begin{framed}
		% \scriptsize
			% \begin{compactdesc}
				% \item[JSP] 
				% % \item[BC] 
				% % \item[]
			% \end{compactdesc}
		% \end{framed}

	% -----
	\subsubsection{Section 39:19} % *DC39-19 
	% -----
		\label{DC39-19}

		Wherefore, go forth, crying with a loud voice, saying: The kingdom of heaven is at hand; crying: Hosanna!  blessed be the name of the Most High God.

		% \begin{framed}
		% \scriptsize
			% \begin{compactdesc}
				% \item[JSP] 
				% % \item[BC] 
				% % \item[]
			% \end{compactdesc}
		% \end{framed}

	% -----
	\subsubsection{Section 39:20} % *DC39-20 
	% -----
		\label{DC39-20}

		Go forth baptizing with water, preparing the way before my face for the time of my coming;

		% \begin{framed}
		% \scriptsize
			% \begin{compactdesc}
				% \item[JSP] 
				% % \item[BC] 
				% % \item[]
			% \end{compactdesc}
		% \end{framed}

	% -----
	\subsubsection{Section 39:21} % *DC39-21 
	% -----
		\label{DC39-21}

		For the time is at hand; the day or the hour no man knoweth; but it surely shall come.

		% \begin{framed}
		% \scriptsize
			% \begin{compactdesc}
				% \item[JSP] 
				% % \item[BC] 
				% % \item[]
			% \end{compactdesc}
		% \end{framed}

	% -----
	\subsubsection{Section 39:22} % *DC39-22 
	% -----
		\label{DC39-22}

		And he that receiveth these things receiveth me; and they shall be gathered unto me in time and in eternity.

		% \begin{framed}
		% \scriptsize
			% \begin{compactdesc}
				% \item[JSP] 
				% % \item[BC] 
				% % \item[]
			% \end{compactdesc}
		% \end{framed}

	% -----
	\subsubsection{Section 39:23} % *DC39-23 
	% -----
		\label{DC39-23}

		And again, it shall come to pass that on as many as ye shall baptize with water, ye shall lay your hands, and they shall receive the gift of the Holy Ghost, and shall be looking forth for the signs of my coming, and shall know me.

		% \begin{framed}
		% \scriptsize
			% \begin{compactdesc}
				% \item[JSP] 
				% % \item[BC] 
				% % \item[]
			% \end{compactdesc}
		% \end{framed}

	% -----
	\subsubsection{Section 39:24} % *DC39-24 
	% -----
		\label{DC39-24}

		Behold, I come quickly.  Even so.  Amen.

		% \begin{framed}
		% \scriptsize
			% \begin{compactdesc}
				% \item[JSP] 
				% % \item[BC] 
				% % \item[]
			% \end{compactdesc}
		% \end{framed}